% !TEX TS-program = pdflatexmk
\documentclass[12pt]{amsart}
\usepackage{multicol}


%\usepackage[parfill]{parskip}    % Activate to begin paragraphs with an empty line rather than an indent

\usepackage[margin=1in]{geometry}


\usepackage{amsmath,amssymb,amsthm,latexsym,graphicx}
\usepackage[normalem]{ulem}
\usepackage{setspace} %used for doublespacing, etc.
\usepackage{hyperref}
\usepackage{cancel}
\usepackage[dvipsnames,usenames]{color}
\usepackage[all]{xy}
\usepackage{fancyhdr}
\pagestyle{fancy}
	\renewcommand{\headrulewidth}{0.5pt} % and the line
	\headsep=1cm
	
\DeclareGraphicsRule{.tif}{png}{.png}{`convert #1 `dirname #1`/`basename #1 .tif`.png}

%Some useful environments.
\newtheorem{theorem}{Theorem}
\newtheorem{corollary}[theorem]{Corollary}
\newtheorem{conjecture}[theorem]{Conjecture}
\newtheorem{lemma}[theorem]{Lemma}
\newtheorem{proposition}[theorem]{Proposition}
\newtheorem{definition}[theorem]{Definition}
\newtheorem{example}[theorem]{Example}
\newtheorem{axiom}{Axiom}
\theoremstyle{remark}
\newtheorem{remark}{Remark}
\newtheorem*{exercise}{Exercise}%[section]

%Some useful shortcuts for our favorite fields
\def\RR{\ensuremath{\mathbb R}} %Note, you can use these WITHOUT entering math mode
\def\NN{\ensuremath{\mathbb N}}
\def\ZZ{\ensuremath{\mathbb Z}}
\def\QQ{{\ensuremath\mathbb Q}}
\def\CC{\ensuremath{\mathbb C}}
\def\EE{{\ensuremath\mathbb E}}

%Some useful shortcuts for formatting lists
\newcommand{\bc}{\begin{center}}
\newcommand{\ec}{\end{center}}
\newcommand{\be}{\begin{enumerate}}
\newcommand{\ee}{\end{enumerate}}
\newcommand{\bi}{\begin{itemize}}
\newcommand{\ei}{\end{itemize}}

%Some useful shortcuts for formatting mathematical symbols
\newcommand{\ol}[1]{\overline{#1}}
\newcommand{\oimp}[1]{\overset{#1}{\iff}} %labeled iff symbol
\newcommand{\bv}[1]{\ensuremath{ \vec{\mathbf{#1}}} } %makes a vector.
\newcommand{\mc}[1]{\ensuremath{\mathcal{#1}}} %put something in caligraphic font
\newcommand{\mpg}[1]{\marginpar{ #1}} %to put comments in margins
\newcommand{\bsl}[1]{\texttt{\symbol{92}{\em #1}}} %for backslashes.
\newcommand{\normale}{\trianglelefteq}
\newcommand{\normal}{\triangleleft}

%Code for formatting the proofs a little nicer for submitted homework
\makeatletter
\renewenvironment{proof}[1][\proofname]{\par\doublespacing
  \pushQED{\qed}%
  \normalfont \topsep6\p@\@plus6\p@\relax
  \list{}{%
    \settowidth{\leftmargin}{\itshape\proofname:\hskip\labelsep}%
    \setlength{\labelwidth}{0pt}%
    \setlength{\itemindent}{-\leftmargin}%
  }%
  \item[\hskip\labelsep\itshape#1\@addpunct{:}]\ignorespaces
}{%
  \popQED\endlist\@endpefalse
  \singlespacing
}
\makeatother


%Commenting tools. You can ignore this stuff unless we're exchanging materials where I'm commenting in detail on how you are writing/using LaTeX.
\usepackage{soul}
\definecolor{highlight}{rgb}{1,0.6,0.6}
\sethlcolor{highlight}
\newcommand{\hlm}[1]{\colorbox{highlight}{$\displaystyle #1$}}
\newtheoremstyle{mycomment}{\topsep}{-0in}{\small \itshape \sffamily}{}{\small \itshape\sffamily}{:}{.5em}{}
\theoremstyle{mycomment}
\newtheorem*{acomment}{\color{BrickRed}{Comment}}
\newcommand{\com}[1]{{\color{OliveGreen}\begin{acomment}{#1} %#2 \color{black} 
\end{acomment}\noindent}}
%\newcommand{\com}[1]{{\color{BrickRed}{\\ Comment:}\color{OliveGreen}{#1} \\}}
\newcommand{\red}[1]{{\color{BrickRed} #1}}
\newcommand{\blue}[1]{{\color{MidnightBlue}#1}}
\newcommand{\green}[1]{{\color{OliveGreen}#1}}
\newcommand{\mwrong}[2]{\red{\cancel{#1}}\green{#2}}
\newcommand{\wrong}[2]{\red{\sout{#1}}\green{#2}}
\definecolor{OliveGreen}{rgb}{.3,.5,.2}
\definecolor{MidnightBlue}{rgb}{.3,.4,.6}
\newcommand{\pts}[1]{\hfill\blue{{#1}/5}}

\chead{MATH 631F}
\pagestyle{fancy}
%Modify these items:
\rhead{\emph{Stefano Fochesatto}}
\lhead{\emph{HW \#1 --- \today}}

\DeclareMathOperator{\lcm}{lcm}
\begin{document}

\thispagestyle{fancy}
\author{Coordinator: Coordinator's Name}


    

\begin{exercise}[0.2.3] Prove that if $n$ is composite then there are integers $a$ and $b$ such that $n$ divides $ab$ but $n$ does not divide either $a$ or $b$.
    \begin{proof} Suppose that $n$ is composite. By definition of composite there exists some positive divisors $a,b \neq 1, n$ such that $ab = n(1)$.
      Clearly $n \mid ab$. Note that $b = n(\frac{1}{a})$ and $a = n(\frac{1}{b})$.
      Since $\frac{1}{a}, \frac{1}{b}$ are $\notin \ZZ$ we have shown that $n \nmid a,b$.
    \end{proof}
    \end{exercise}

\begin{exercise}[0.2.10] Prove for any given positive integer $N$ there exist only finitely many integers $n$ 
    with $\varphi(n) = N$ where $\varphi$ denotes Euler's $\varphi$-function. Conclude in particular that 
    $\varphi(n)$ tends to infinity as $n$ tends to infinity.
    \begin{proof} 
        Let $N \in \ZZ^{+}$, and $X = \{n \in \ZZ: \varphi(n) = N\}$. By the Fundamental Theorem of Arithmetic, $n = p_1^{\alpha_1}p_2^{\alpha_2} \dots p_s^{\alpha_s}$
        for primes $p_i$ and $\alpha_i \in \ZZ^+$. Consider prime $p$, the largest of primes $p_i$. It follows that, 
        \begin{equation*}
        N = \varphi(n) = \prod_{i = 1}^s p_i^{\alpha_i - 1}(p_i - 1) \geq (p - 1).
        \end{equation*}
        Therefore for all follows that $p_i \leq p \leq N+1$.\\

        Note that
        \begin{equation*}
        N = \varphi(n) = \prod_{i = 1}^s p_i^{\alpha_i - 1}(p_i - 1) \geq \prod_{i = 1}^s p_i^{\alpha_i - 1} \geq \prod_{i = 1}^s 2^{\alpha_i - 1}. 
        \end{equation*}
        Consider $\alpha$, the largest element in the set of $\alpha_i$. It follows that, 
        \begin{equation*}
        N \geq \prod_{i = 1}^s 2^{\alpha_i - 1} \geq 2^{\alpha - 1}.
        \end{equation*}
        Therefore $\alpha_1 \leq \alpha \leq \lceil\log_2(N) + 1\rceil$.
        Since $p_i$ is finite, and $s_i$ is finite, and whose elements are bounded above by some function of $N$ then the set,
        \begin{equation*}
        A  = \{x = \prod_{i = 1}^s q_{i}^{\beta_i}: q_i \in p_i, \beta_i \in \alpha_i\}
        \end{equation*} 
        is finite with $a \in A$ bounded above by some function of $N$. 
        Note that $X$ is a subset of $A$ and therefore $X$ is finite, and for all $n \in X$, $n$ is bounded above by some function of $N$. Thus as $n$ tends to infinity $\varphi(n)$ also tends to infinity.  
    \end{proof}
\end{exercise}

\begin{exercise}[0.3.9]Prove that the square of any odd integer always leaves a remainder of 1 when divided by 8.
    \begin{proof} Suppose $n \in \ZZ$ is odd. By the definition, for some $i \in \ZZ$, $n = 2(i) + 1$.
      Consider $n^2$,
      \begin{equation*}
        n^2 = (2(i) + 1)(2(i) + 1) = 4(i)^2 + 4(i) + 1 = 4i(i + 1) + 1.
      \end{equation*}
      Note that when $i$ is odd, $i + 1$ must be even and vice versa. Therefore we can always factor out a $2$ from the product $i(i + 1)$ giving us
      for some $j \in \ZZ$, 
      \begin{equation*}s
        n^2 = 8(j) + 1.
      \end{equation*}
    \end{proof}
  \end{exercise}
  
    


\begin{exercise}[1.1.25] Prove that if $x^2 = \epsilon$ for all $x \in G$ then $G$ is abelian. 
    \begin{proof} Suppose $x^2 = \epsilon$ for all $x \in G$. With some algebra we get, 
      \begin{align*}
        x^2 &= \epsilon,\\
        x^2(x^{-1}) &= \epsilon(x^{-1}),\\
        x&= x^{-1}.
      \end{align*}
      Now consider some $a, b \in G$, and with (4) of Proposition 1 we can consider the following, 
      \begin{equation*}
        ab = (ab)^{-1} = b^{-1}a^{-1} = ba.  
      \end{equation*}
      Thus $G$ is abelian. 
    \end{proof}
  \end{exercise}
      


\begin{exercise}[1.3.4b] Compute the order of the elements in $S_4$.
    \begin{proof}[Solution:] First recall from the end of Section 1.3 
      that that the order of a permutation is the l.c.m of the lengths of the cycles in it's cycle decomposition. Now we can compute the order of each permutation in $S_4$,
      \begin{center}
      \begin{minipage}{.45\textwidth}
        \begin{tabular}{c | c } 
          Permutation $(\sigma)$ & $|\sigma|$ \\ 
          \hline
          (1) & 1\\
          (12) & 2\\
          (13) & 2\\
          (14) & 2\\
          (23) & 2\\
          (24) & 2\\
          (34) & 2
        \end{tabular}
      \end{minipage}
      \hfill
      \begin{minipage}{.45\textwidth}
        \begin{tabular}{c | c } 
          Permutation $(\sigma)$ & $|\sigma|$ \\ 
          \hline
          (123) & 3\\
          (124) & 3\\
          (132) & 3\\
          (134) & 3\\
          (142) & 3\\
          (143) & 3\\
          (234) & 3\\
          (243) & 3
        \end{tabular}
      \end{minipage}     
    \end{center}
  
    \begin{center}
      \begin{minipage}{.45\textwidth}
            \begin{tabular}{c | c } 
              Permutation $(\sigma)$ & $|\sigma|$ \\ 
              \hline
              (1234) & 4\\
              (1243) & 4\\
              (1324) & 4\\
              (1342) & 4\\
              (1423) & 4\\
              (1432) & 4
            \end{tabular}
          \end{minipage}
          \hfill
          \begin{minipage}{.45\textwidth}
              \begin{tabular}{c | c } 
                Permutation $(\sigma)$ & $|\sigma|$ \\ 
                \hline
                (12)(34)& 2\\
                (13)(24)& 2\\
                (14)(23)& 2
              \end{tabular}
            \end{minipage}     
          \end{center}
          
    \end{proof}
    
  \end{exercise}


  \end{document}